Aus $pV = nRT$ folgt $V = nR\cdot\frac{T}{p}$, in den Einheiten des Diagramms $V/V_0 = \frac{T/T_0}{p/p_0}$. Im $p(T)$-Diagramm ist eine Isobare waagrecht, eine Isotherme senkrecht und eine Isochore eine Gerade durch den Ursprung ($p \propto T$).
\begin{abcliste}
\abc A $\to$ B ist waagrecht: $\result{\text{isobar}}$. B $\to$ C ist senkrecht: $\result{\text{isotherm}}$. C $\to$ D liegt auf einer Geraden durch den Ursprung: $\result{\text{isochor}}$.

\abc Die Volumen der Zustände:
\begin{center}
\begin{tabular}{lcccc}
 & A & B & C & D \\ \hline
 $T$ in $T_0$ & 4 & 2 & 2 & 6 \\
 $p$ in $p_0$ & 1 & 1 & 2 & 6 \\
 $V$ in $V_0$ & 4 & 2 & 1 & 1
\end{tabular}
\end{center}
A $\to$ B: Bei gleichem Druck halbiert sich $T$, also auch $V$: $\result{V \text{ halbiert sich}}$ (von $4\,V_0$ auf $2\,V_0$). B $\to$ C: Bei gleicher Temperatur verdoppelt sich $p$, also $\result{V \text{ halbiert sich}}$ (auf $V_0$). C $\to$ D: isochor, $\result{V \text{ bleibt gleich}}$ ($V_0$), während sich $p$ und $T$ verdreifachen.

\abc Entlang D $\to$ A sinken sowohl $p$ als auch $T$, aber nicht im selben Verhältnis: $T/p$ ändert sich, also auch $V$.
\begin{align}
 \text{D: } \frac{V}{V_0} &= \frac{6}{6} = 1 & \text{A: } \frac{V}{V_0} &= \frac{4}{1} = 4
\end{align}
Der Druck fällt auf einen Sechstel, die Temperatur nur auf zwei Drittel. Das Gas $\result{\text{dehnt sich aus}}$, sein Volumen $\result{\text{vervierfacht sich}}$. Im Diagramm: Je flacher die Verbindung eines Punktes mit dem Ursprung ist, desto grösser ist $V$; die Gerade OD hat die Steigung 1, die Gerade OA nur $\frac14$.
\end{abcliste}
